\documentclass[../cours_1re.tex]{subfiles}
\begin{document}
\setcounter{chapter}{23}
\chapter{Algorithmique, listes et calcul fiable}
\label{ch:24}
\begin{objectif}
Écrire, lire et vérifier des programmes courts. Prérequis : variables, tests, boucles et fonctions de seconde.
Les fichiers exécutables et leurs vérifications se trouvent dans \texttt{premiere/python/}.
\end{objectif}
\section{Données, fonctions et contrats}
Les types usuels sont \texttt{int}, \texttt{float}, \texttt{bool} et \texttt{str}.
L'affectation \texttt{x = x + 1} modifie une valeur ; ce n'est pas une équation.
Une fonction documente ses entrées, sa sortie et ses conditions :
par exemple la factorielle exige un entier $n\ge0$.
\begin{lstlisting}[language=Python]
def factorielle(n):
    if type(n) is not int or n < 0:
        raise ValueError("entier positif ou nul attendu")
    produit = 1
    for k in range(1, n + 1):
        produit *= k
    return produit
\end{lstlisting}
Après l'étape $k$, \texttt{produit} vaut $k!$ : c'est un invariant.
L'initialisation correspond à $0!=1$ ; la multiplication maintient l'invariant.
La boucle possède un nombre fini d'étapes, donc termine.
Un test sur $5!$ détecte certaines erreurs ; l'invariant prouve tous les cas admissibles.

\section{Listes : ordre, indices et modifications}
\begin{lstlisting}[language=Python]
termes = [2, 5, 8]             # liste en extension
termes.append(11)              # ajoute a la fin
premier = termes[0]            # indices a partir de zero
dernier = termes[-1]
termes[1] = 6                 # remplace un element
del termes[1]                 # supprime l'element d'indice 1
retire = termes.pop()         # retire et renvoie le dernier
carres = [k*k for k in range(6)]
somme = 0
for valeur in carres:
    somme += valeur
\end{lstlisting}
Après ces instructions, \texttt{termes} vaut \texttt{[2, 8]} et
\texttt{carres} vaut \texttt{[0, 1, 4, 9, 16, 25]}. Une liste accepte
les répétitions et conserve un ordre. Une liste de $n$ éléments a ses indices
positifs de $0$ à $n-1$ ; l'indice $n$ provoque une erreur.
\texttt{b = a} désigne la même liste que \texttt{a}.
Pour une liste simple de nombres, \texttt{b = a.copy()} crée une copie indépendante.
Une copie superficielle ne copie pas récursivement les sous-listes.

\section{Récurrences, seuils et arrêt}
Pour Fibonacci, initialiser $(a,b)=(0,1)$, enregistrer $a$, puis faire
\texttt{a, b = b, a + b}. Avant l'itération $k$, les valeurs sont
$(F_k,F_{k+1})$. L'affectation simultanée évite de perdre l'ancienne valeur.
Pour une suite à seuil, préciser si l'on cherche $u_n>T$ ou $u_n\ge T$.
Si l'existence du seuil n'est pas prouvée, fixer un nombre maximal d'itérations
et signaler explicitement son dépassement.

La suite de Syracuse part d'un entier positif : on divise un terme pair
par $2$ et on remplace un terme impair $u$ par $3u+1$.
Le programme joint s'arrête à $1$ ou à une borne choisie.
Il ne suppose pas résolue la conjecture d'arrêt pour tout entier initial.
Le calcul entier \texttt{u // 2} évite ici une conversion inutile en flottant.

\section{Les algorithmes d'analyse et de géométrie}
Pour approcher une dérivée en $a$, construire la liste des
$(f(a+h)-f(a))/h$ pour plusieurs $h\ne0$. Pour $f(x)=x^2$,
la valeur exacte est $2a+h$ : cet exemple permet de contrôler le code.
Des $h$ extrêmement petits peuvent aggraver l'erreur d'arrondi par soustraction.

Euler met à jour \texttt{y = y + h * f(t, y)}, puis le temps.
Newton met à jour \texttt{x = x - f(x) / df(x)} seulement si le dénominateur
n'est pas nul. Une petite différence entre itérés ne garantit une erreur
petite que si une estimation mathématique la justifie (chapitre 29).

Pour Archimède, partir de l'hexagone inscrit au cercle unité : $n=6$, $s=1$.
Les demi-périmètres inscrit et circonscrit encadrent $\pi$ :
\[
L=\frac{ns}{2}<\pi<\frac{ns}{2\sqrt{1-s^2/4}}=U.
\]
Le passage au double du nombre de côtés utilise
$s_{\rm nouveau}=s/\sqrt{2+\sqrt{4-s^2}}$ (chapitre 14).
Cette forme évite la soustraction de nombres voisins.
L'encadrement est exact en arithmétique réelle ; des flottants ordinaires
ne constituent pas un encadrement certifié pour un nombre arbitraire d'itérations.

\section{Simuler et traiter des données}
Pour une loi finie de probabilités $(p_1,\ldots,p_r)$, cumuler ces probabilités
et comparer un uniforme $U$ aux cumuls : le $i$-ième intervalle a une longueur
$p_i$, donc sélectionne la valeur $x_i$ avec cette probabilité dans le modèle.
Utiliser ensuite des compréhensions de listes pour former des échantillons,
des moyennes et la proportion demandée au chapitre 23.

Pour étudier la fréquence des lettres, le module normalise les accents,
retient les lettres de \texttt{a} à \texttt{z}, puis utilise une liste de
26 effectifs. Diviser par le nombre de lettres retenues donne les fréquences.
Deux courtes phrases françaises et anglaises sont fournies comme données
d'exercice. Elles permettent de vérifier le programme ; elles ne suffisent
pas à établir la distribution des lettres d'une langue.

Dans un tableur, une suite $u_{n+1}=0{,}8u_n+30$ peut être placée en colonne B :
indice $0$ en A2, $u_0$ en B2, puis \texttt{=A2+1} en A3 et
\texttt{=0.8*B2+30} en B3 (adapter la virgule décimale à la configuration).
Recopier vers le bas ; un paramètre fixe placé en D1 s'utilise avec
\texttt{\$D\$1}. Pour une suite, choisir un nuage de points : l'indice
est une abscisse numérique, pas une catégorie arbitrairement espacée.

\begin{finchapitre}
\section{Tester et discuter la fiabilité}
Le script de vérification contrôle des valeurs exactes, les entrées invalides,
des identités de récurrence, les bornes de Newton et l'évolution de l'erreur
d'Euler. Les simulations ont des graines fixes et sont comparées à des bornes
raisonnables, sans exiger une fréquence théorique exacte.
Séparer erreur de modèle, erreur de discrétisation, fluctuation aléatoire
et erreur d'arrondi aide à corriger la bonne cause.


\section{Utilité et prolongements}
Un algorithme rend une définition calculable. Le contrat, l'invariant et
la preuve d'arrêt en établissent la signification ; les tests contrôlent l'implémentation.

En licence puis en master, analyse numérique, probabilités et optimisation
associeront systématiquement méthode de calcul, démonstration de convergence
et estimation de l'erreur. Les listes préparent les vecteurs de données.
\begin{retenir}
Préciser les indices, vérifier les conditions et borner les boucles incertaines.
Calculer plus longtemps ou avec des nombres plus petits ne garantit pas une meilleure précision.
\end{retenir}
\end{finchapitre}
\end{document}
